\begin{abstract}
Graphs with millions of nodes are now routine in recommendation, biology, and the analysis of social networks, yet most spectral learning methods still assume that a full eigendecomposition of the graph Laplacian is affordable. This thesis studies what can be learned from a small number of approximate eigenvectors. We give error bounds for spectral embeddings computed with randomized sketches, propose a filtering scheme that needs only sparse matrix--vector products, and show on six benchmarks that the resulting models match dense baselines at a fraction of the memory.
\end{abstract}
